\documentclass{article}
\usepackage[margin=0.75in]{geometry}
\usepackage{xcolor}
\usepackage{amsmath}
\usepackage{amssymb}
\usepackage{fontspec}
\usepackage{tikz}
\usetikzlibrary{tikzmark}
\usepackage{enumitem}
\usepackage{multicol}

\definecolor{bglight}{HTML}{FFFFFF}
\definecolor{textlight}{HTML}{000000}
\definecolor{subtext}{HTML}{666666}

\pagecolor{bglight}
\color{textlight}

\pagestyle{empty}

% #1 = number, #2 = node name used later for the connector lines
\newcommand{\stepnum}[2]{%
  \tikzmarknode{#2}{%
    \begin{tikzpicture}[baseline=(char.base)]
      \node[draw=textlight, circle, inner sep=1.5pt, font=\sffamily\small] (char) {#1};
    \end{tikzpicture}%
  }%
}

\begin{document}
\sffamily

% ================= HEADER =================
\begin{center}
  {\Huge \bfseries Making 10s Exercise Sheet} \\[0.4em]
\end{center}
\vspace{1em}

\section*{Part 1: Addition}
% ================= EXAMPLE SECTION =================
{\Large \textbf{Example: $11 + 51 + 101$}}
\vspace{0.5em}

\begin{description}[leftmargin=2em]

  \item[\stepnum{1}{s1}] \textbf{Decompose each number into tens and ones}\\[0.2em]
  Separating values makes mental math simple
  \begin{itemize}
    \item $11  = 10  + 1$
    \item $51  = 50  + 1$
    \item $101 = 100 + 1$
  \end{itemize}
  \vspace{0.6em}

  \item[\stepnum{2}{s2}] \textbf{Group and add all the tens together}\\[0.4em]
  Combine the multiples of 10:
  \[
    10 + 50 + 100 = 160
  \]
  \vspace{0.6em}

  \item[\stepnum{3}{s3}] \textbf{Group and add all the remaining ones}\\[0.4em]
  Combine the ones:
  \[
    \mathbf{1} + \mathbf{1} + \mathbf{1} = \mathbf{3}
  \]
  \vspace{0.6em}

  \item[\stepnum{4}{s4}] \textbf{Add the totals together}\\[0.4em]
  Combine the sum of the tens with the sum of the ones:
  \[
    160 + \mathbf{3} = \mathbf{163}
  \]

\end{description}

\begin{tikzpicture}[remember picture, overlay]
  \draw[dotted, thick, textlight!60] (s1.south) -- (s2.north);
  \draw[dotted, thick, textlight!60] (s2.south) -- (s3.north);
  \draw[dotted, thick, textlight!60] (s3.south) -- (s4.north);
\end{tikzpicture}

\vspace{1em}
\noindent\textcolor{subtext!40}{\rule{\linewidth}{0.5pt}}
\vspace{1.5em}

% ================= QUESTIONS SECTION =================
\begingroup
\parindent=0pt
{\Large \textbf{Practice Questions}}\par\vspace{0.4em}
\endgroup
\vspace{1em}

\begin{multicols}{2}
\begin{enumerate}[label=\textbf{\arabic*.}, itemsep=1.8em, leftmargin=2em]
  \item $11 + 31 + 61 = {\color{subtext}10 + 30 + 60 + 3 = 103}$
  \item $21 + 41 + 51 = \underline{\hspace{2.5cm}}$
  \item $11 + 21 + 31 + 51 = \underline{\hspace{2.5cm}}$
  \item $31 + 51 + 71 = \underline{\hspace{2.5cm}}$
  \item $11 + 41 + 81 = \underline{\hspace{2.5cm}}$
  \item $21 + 31 + 41 + 51 = \underline{\hspace{2.5cm}}$
  \item $11 + 61 + 91 = \underline{\hspace{2.5cm}}$
  \item $41 + 51 + 61 = \underline{\hspace{2.5cm}}$
  \item $11 + 11 + 31 + 41 = \underline{\hspace{2.5cm}}$
  \item $21 + 51 + 81 = \underline{\hspace{2.5cm}}$
\end{enumerate}
\end{multicols}

\section*{Part 2: Subtraction}
% ================= EXAMPLE SECTION =================
{\Large \textbf{Example: $110 - 21 - 31$}}
\vspace{0.5em}

\begin{description}[leftmargin=2em]

  \item[\stepnum{1}{s1}] \textbf{Decompose the numbers to be subtracted}\\[0.2em]
  Break the numbers into tens and ones
  \begin{itemize}
    \item $21  = 20  + 1$
    \item $31  = 30  + 1$
  \end{itemize}
  \vspace{0.6em}

  \item[\stepnum{2}{s2}] \textbf{Subtract the tens}\\[0.4em]
  Subtract the multiples of 10 from the starting number ($110$):
  \[
    110 - 20 - 30 = 60
  \]
  \vspace{0.6em}

  \item[\stepnum{3}{s3}] \textbf{Subtract the remaining ones from the result}\\[0.4em]
  Take the remaining ones ($1$ and $1$) away from $60$:
  \[
    60 - \mathbf{1} - \mathbf{1} = \mathbf{58}
  \]

\end{description}

\begin{tikzpicture}[remember picture, overlay]
  \draw[dotted, thick, textlight!60] (s1.south) -- (s2.north);
  \draw[dotted, thick, textlight!60] (s2.south) -- (s3.north);
\end{tikzpicture}

\vspace{1em}
\noindent\textcolor{subtext!40}{\rule{\linewidth}{0.5pt}}
\vspace{1.5em}

% ================= QUESTIONS SECTION =================
\begingroup
\parindent=0pt

{\Large \textbf{Practice Questions}}\par\vspace{0.4em}
\endgroup
\vspace{1em}

\begin{multicols}{2}
\begin{enumerate}[label=\textbf{\arabic*.}, itemsep=1.8em, leftmargin=2em]
  \item $\begin{aligned}[t] 120 - 21 - 31 &= {\color{subtext}120 - 20 - 30 - 1 - 1} \\ &= {\color{subtext}70 - 1 - 1 = 68} \\ &\vphantom{70 - 1 - 1 = 68}\end{aligned}$
  \item $\begin{aligned}[t] 100 - 11 - 21 &= \underline{\hspace{2.5cm}} \\ &= \underline{\hspace{2.5cm}} \\ &\vphantom{70 - 1 - 1 = 68}\end{aligned}$
  \item $\begin{aligned}[t] 150 - 41 - 51 &= \underline{\hspace{2.5cm}} \\ &= \underline{\hspace{2.5cm}} \\ &\vphantom{70 - 1 - 1 = 68}\end{aligned}$
  \item $\begin{aligned}[t] 180 - 31 - 41 &= \underline{\hspace{2.5cm}} \\ &= \underline{\hspace{2.5cm}} \\ &\vphantom{70 - 1 - 1 = 68}\end{aligned}$
  \item $\begin{aligned}[t] 110 - 11 - 51 &= \underline{\hspace{2.5cm}} \\ &= \underline{\hspace{2.5cm}} \\ &\vphantom{70 - 1 - 1 = 68}\end{aligned}$
  \item $\begin{aligned}[t] 130 - 31 - 21 - 11 &= \underline{\hspace{2.5cm}} \\ &= \underline{\hspace{2.5cm}} \\ &\vphantom{70 - 1 - 1 = 68}\end{aligned}$
  \item $\begin{aligned}[t] 160 - 51 - 61 &= \underline{\hspace{2.5cm}} \\ &= \underline{\hspace{2.5cm}} \\ &\vphantom{70 - 1 - 1 = 68}\end{aligned}$
  \item $\begin{aligned}[t] 140 - 21 - 41 &= \underline{\hspace{2.5cm}} \\ &= \underline{\hspace{2.5cm}} \\ &\vphantom{70 - 1 - 1 = 68}\end{aligned}$
  \item $\begin{aligned}[t] 200 - 31 - 51 - 61 &= \underline{\hspace{2.5cm}} \\ &= \underline{\hspace{2.5cm}} \\ &\vphantom{70 - 1 - 1 = 68}\end{aligned}$
  \item $\begin{aligned}[t] 170 - 41 - 31 - 21 &= \underline{\hspace{2.5cm}} \\ &= \underline{\hspace{2.5cm}} \\ &\vphantom{70 - 1 - 1 = 68}\end{aligned}$
\end{enumerate}
\end{multicols}
\end{document}